Orbiter can be used to create raytracing 3D-graphics.
Orbiter serves as a front end for the raytracing software Povray~\cite{povray}.
This is a multi step process: 
A 3D scene is defined through orbiter commands. 
Next, Orbiter produces Povray files. 
After that, the 
povray files are processed through povray, and turned into graphics files (png), called frames. 
The frames can be turned into a video 
by using tools like ffmpeg (see Section~\ref{sec:animations}).
By default, an rotational animation is produced. 
The main commands are listed in Table~\ref{tab:interface:povray}.
\orbitertableofcommands{interface_povray}

\bigskip

To describe a povray job, the commands in Table~\ref{tab:povray:job} can be used.
\orbitertableofcommands{povray_job_description}

\bigskip

The Orbiter Povray interface requires some general information about the animation, 
the camera position, the boundary box for clipping, 
the font size for text and others.
Tables~\ref{tab:povray1}-~\ref{tab:povray2} 
list the commands to control the 3D-povray frontend.
\orbitertableofcommands{video_draw_options_1}
\orbitertableofcommands{video_draw_options_2}

\bigskip


A povray graphics starts with a description of the scene.
The scene will be rendered to the final frame. 
A scene is composed of geometric objects. 
Various types of objects are available: 
points, lines, planes, faces, algebraic surfaces, reguli, 3D-text,  
and others.
Some complex objects are predefined, 
for instance the Eckardt surface. 
Once the objects are defined, output commands can be invoked 
to draw them in various colors and with various options.
At times, there are many objects in one scene. 
In order to ease drawing, objects can be grouped together. 
Only objects of the same type may be grouped. 
A group of object can be drawn at once.
Tables~\ref{tab:scene1} and~\ref{tab:scene2} 
summarize the Orbiter commands to build objects of a 3D scene.
\orbitertableofcommands{scene_1}
\orbitertableofcommands{scene_2}
The commands in Table~\ref{tab:3doutput} are used to build the scene.
Each of these commands applies to a group of objects of the same kind. 
Groups of objects are created using the commands in 
Table~\ref{tab:scene2} which start with \verb'group_of'.
\orbitertableofcommands{drawable_set_of_objects}
Here is a simple example which combines scene building and graphical output.
The example creates a cube with vertices, edges and faces:

\sourcecode{cube}

This command instructs Orbiter to create 30 povray files (extension .pov), 
one for each frame of a rotating scene. 
The scene containes a cube whose vertices are shown in chrome, 
whose edges are in red, and whose faces are yellow and transparent. 
The cube turns around a vertical axis of symmetry.
Here is the first frame of the result:
\orbiteroutput{GRAPHICS/WRAPPERS/povray_cube}
The coordinates of the cube are stored in an object file \verb'cube_centered.obj'.
The content of this file is:
\orbiteroutput{GRAPHICS/WRAPPERS/povray_cube_obj_file}
The monkey saddle is a cubic surface, given by the equation 
$$
z = x^3-3xy^2
$$
The next example plots the surface knowns as the monkey saddle. 
The tangent plane at $(0,0,0)$ is drawn as well. 
An animation is created by rotating the scene around the $z$-axis.

\sourcecodevariable{MONKEY_SADDLE_CUBIC}


\sourcecode{monkey}

Here is one of the frames that are created:
\orbiteroutput{GRAPHICS/WRAPPERS/povray_monkey}
The Eckardt surface is given by the equation
$$
\frac{5}{2}xyz-(x^2+y^2+z^2)+1=0.
$$
We use the following code to plot the surface and the lines on it. 
The Schl\"afli labeling of the lines is indicated.

\sourcecode{Eckardt}

The rendering yields:
\orbiteroutput{GRAPHICS/WRAPPERS/povray_Eckardt}






\bigskip

The Endrass octic~\cite{Endrass97} is the algebraic surface given by the equation

\input GRAPHICS/WRAPPERS/endrass_octic_equation.tex

The following Orbiter command creates a povray graphics of the octic. 

\sourcecodevariable{ENDRASS_OCTIC_LEX_165}


\sourcecode{endrass8}

Here is the rendering:
\orbiteroutput{GRAPHICS/WRAPPERS/Endrass_octic}
It includes coordinate axes and the $x,y$-plane. 

\bigskip

